\begin{afterword}
\summaryhead

The post follows ``The Second Law of Thermodynamics, and Engines of Cognition,'' which concluded that forming accurate beliefs requires observation. Its moral is that the laws of probability are as binding as the laws of physics. People taught that certain beliefs are orders tend to treat a probable belief as a suggestion, as when a lottery buyer says that no one can prove the ticket will lose. But no one expects a smashed egg to reassemble or boiling water to cool a hand, though both are only improbable.

The author recalls a man whose spreadsheet showed that his wheels and gears produced reactionless thrust, which classical mechanics rules out, so the spreadsheet had to contain a mistake. People who believe without evidence, the post says, build similar structures of argument, and somewhere in them a tiny probability is turned into a large one. Each added step, like each added gear, only lowers the probability of the whole.

\responsehead

The physics in the post is right, and so is its main point about low probabilities. The second law holds with overwhelming probability, not certainty; a correct calculation in classical mechanics cannot show reactionless thrust; guessing the exact state of boiling water is no more likely than being spared by chance. A tiny probability is a belief one must hold, not a permission to believe otherwise.

The trouble is in the step from these cases to the post's target, beliefs held without evidence. In every case the post works through, both sides agree that the probability is tiny: physics fixes it for the egg and the water, and the rules of the draw fix it for the lottery. The laws of probability then decide what follows. For a belief held without evidence, the tiny starting probability is the point in question, and the laws of probability do not supply it. They govern how probabilities combine and change, and a subjective Bayesian may hold any coherent prior. The post's reason for a tiny start is one sentence, that possibility holds ``far more thin air than ground.'' That is a claim about priors. The post presents it beside laws it calls ``harder than steel'' without saying that it is of a different kind.

The argument the post warns against is not shown. It says that such reasoners ``often build vast edifices of justification'' and that there is ``always some step'' where a tiny probability becomes a large one. No such argument is examined. The one real case, a physics spreadsheet with a calculation error, is offered as an analogy for such arguments, not as one of them.

The post also leans on the previous one. Its account of heat as energy that cannot be extracted because its state is unknown is Jaynes's disputed view, and the conclusion it quotes includes the claim that observing costs thermodynamic work, which the notes on that post discuss.

In short: correct physics and a sound point about low probabilities, carried over to beliefs without evidence by an analogy that assumes the tiny starting probability it needs, with no example of the arguments it describes.
\end{afterword}
